\chapter{面向服务应用的框架}
\label{chap:soa-ch}

\section{本章目的}

本章不是要对所有服务框架做宽泛综述，而是回顾已有系统已经解决了什么，并
说明为什么在这些系统存在的情况下仍然需要 NDNSF。关键问题不是 gRPC、ROS、
微服务或 NDN Service Call 是否有用；它们显然有用。关键问题是：为什么动
态、安全、协作型应用仍然需要另一个框架。

\section{主机中心 RPC 和微服务框架}

RPC 系统（如 gRPC~\cite{grpc}）、HTTP/REST 系统~\cite{fielding2000rest}
以及微服务平台已经成为构建分布式软件的主流方式。它们提供类型化 API、传
输绑定、负载均衡、服务注册和生产环境工具。当服务能够由稳定端点标识、信
道能够建立并保持、安全性能够表示为信道安全加服务端授权时，这些系统工作
得很好。

然而，在移动和间歇连接环境下，这种抽象不再自然。UAV 视频服务用户关心
的是来自无人机 A 的最新帧，而不是无人机 A 当前通过哪个 IP 地址可达。分
布式推理客户端关心的是某个模型阶段的服务名，而不是固定 GPU 服务器。任
务规划应用可能需要任意满足角色要求的无人机，而不是某个预配置端点。主机
中心框架可以支持这些场景，但发现、重试、状态跟踪和应用特定协调的负担会
转移到框架之外。

\subsection{这些系统已经解决的问题}

RPC 和微服务系统解决了 NDNSF 不应忽视的若干问题。首先，它们提供熟悉的
编程模型：服务有接口、请求类型、响应类型、错误处理和 deadline。其次，
围绕 gRPC 和 REST 的生产系统提供成熟的负载均衡、日志、监控、部署和版本
管理工具。第三，TLS 等信道安全机制提供了保护客户端服务器通信的实用方式。
对于稳定云服务，这些能力通常已经足够。

NDNSF 并不试图否定这些工程经验，而是关注当服务提供者移动、间歇连接、复
制或嵌入资源受限边缘环境时，哪些假设会失效。在这些场景中，一次服务调用
不仅是信道操作，也是关于哪个提供者当前被授权、可达、有能力且值得选择的
决策。现有框架可以实现这些逻辑，但通常需要分散在服务发现系统、负载均衡
器、应用重试、中间件插件和服务特定策略代码中。

\subsection{NDNSF 需要解决的开放问题}

第一个问题是端点依赖。即使现代 service mesh 用逻辑服务名隐藏 IP 地址，
底层通信仍然依赖主机中心路由到达某个端点。第二个问题是数据生命周期。一
旦响应、视频 chunk、模型分片或遥测日志被存储到原信道之外，应用必须额外
定义验证和加密模型。第三个问题是服务提供者多样性。RPC 通常在请求发送前
就选择单个服务器；NDNSF 则允许多个提供者观察请求、决定是否 ACK，并暴露
帮助用户选择的元数据。这些差异促使我们设计数据中心服务框架，而不仅仅是
另一个 IP endpoint wrapper。

\section{机器人中间件和服务组合}

ROS 和 ROS2~\cite{ros,ros2} 等机器人中间件为机器人组件提供了更贴近应
用的模型。节点可以发布 topic、调用 service，并通过标准消息类型交互。与
早期 ROS 相比，ROS2 改进了 QoS 控制和部署灵活性。这些框架对 UAV 应用
非常重要，因为它们将传感器、控制、遥测和感知视为可组合模块。

剩余缺口是：机器人中间件不是通用的数据中心网络安全框架。Trust anchor、
证书部署、数据静态保护、多提供者选择和跨域服务授权并不是其核心抽象。此
外，部署在地面站、多架无人机和间歇无线链路上的 UAV 应用仍然需要在单一
本地机器人域之外可靠工作的命名、验证和恢复机制。

\subsection{对 UAV 服务容器的启示}

ROS 对本论文特别重要，因为它表明机器人应用天然是面向服务的。UAV 系统
不是单一通信流，而是由飞控命令、遥测流、摄像头帧、地图、任务规划、目标
检测、日志和健康监测共同组成。ROS 式组合验证了 UAV 软件应该被分解成服
务的思想。

然而，NDNSF 面向的是不同层次。ROS 和 ROS2 是本地及分布式机器人组件的
强中间件，而 NDNSF 关注基于 NDN 的命名、安全、网络可见服务调用。在拟
实现的 UAV 应用中，无人机 app 仍可被视为服务容器，但地面站命令、遥测发
现、视频获取和跨无人机协作等外部交互使用 NDNSF 服务名和数据中心验证。
因此，NDNSF 不是飞控或机器人中间件的替代品，而是连接这些组件并使它们能
在动态 NDN 环境中运行的网络服务框架。

\section{消息队列和事件系统}

消息队列和发布订阅系统在时间上解耦生产者和消费者~\cite{messagequeue}。
它们适合遥测、日志、异步工作流和弹性后端处理。这与 NDNSF 相关，因为许
多服务应用同时包含请求响应动作和事件流。例如，无人机可能以服务请求形式
接收命令，同时异步发布遥测或视频元数据；分布式推理流水线也可能在处理用
户请求的同时发布中间结果和执行事件。

然而，许多消息队列部署依赖 broker 或逻辑中心化基础设施。安全性通常通过
broker 访问控制和信道加密执行，而不是由可在任何存储或转发位置验证的数
据对象来保证。对于移动或断连边缘环境，broker 的可用性和部署位置可能成为
限制。此外，队列本身并不定义用户如何在多个服务提供者中选择、如何验证提
供者是否被授权提供特定服务，或如何以可验证数据形式获取大型服务 artifact。

\section{NDN 计算和服务框架}

NDN 已经为面向服务系统提供了有用原语：层次化名字、接收端驱动获取、网内
缓存、自适应转发和 Data 签名~\cite{ndn}。NDN 上的服务框架必须弥合这些
原语和应用层服务语义之间的差距。已有 NDN 工作探索了该设计空间的多个区
域。

\subsection{函数中心计算}

Named Function Networking（NFN）将计算视为一等命名对象，并将函数表达式
嵌入名字~\cite{nfn}。这种模型很强大，因为它允许计算和数据获取在命名系
统内组合，也展示了内容名和计算名深度集成的可能性。

代价是复杂性。函数中心命名具有很强表达能力，但执行语义、信任、缓存和部
署也更难推理。UAV 控制、实时视频和分布式 AI 推理需要显式服务授权、服务
提供者选择、长时间状态和应用定义的执行策略。因此，NDNSF 选择服务调用抽
象，而不是完全通用的函数表达式抽象。

\subsection{Serverless 和边缘型 NDN 框架}

NFaaS 将 NDN 推向 serverless computing，把函数视为可在边缘位置执行的可
部署单元~\cite{nfaas}。其他边缘框架，例如面向信息中心 IoV 的 CFEC，也探
索低延迟微服务执行和动态放置~\cite{cfec}。这些系统突出一个重要方向：服
务执行应靠近用户、数据源或资源丰富的边缘节点。

开放问题在于，动态放置只是服务应用支持的一部分。应用还需要稳定 API、用
户和提供者授权、多提供者选择、大对象处理以及清晰部署流程。NDNSF 关注这
些运行时服务语义。它也可以与未来放置机制共存；例如，放置算法可以决定哪
些无人机、存储组件或推理提供者应发布某个 NDNSF 服务。

\subsection{远程调用框架}

RICE 在 ICN 中提供远程方法调用，使用 handshake 和 thunk 表示计算参数与
结果~\cite{rice}。NSC 为 NDN 提出 Named Service Calls，并提供类似 RPC
的远程执行抽象~\cite{nsc}。这些系统与 NDNSF 目标接近，因为它们直接处理
服务调用，而不只是内容获取。

主要限制是它们没有提供本论文目标应用所需的完整机制。NSC 主要是一对一：
调用者联系 executor，协议不是围绕同时服务提供者发现、ACK 元数据和多提供
者选择设计的。RICE 提供更强调用语义，但需要修改 forwarding plane。NDNSF
将框架保留在应用/运行时层，并使用 NDN Data 发布和 Sync 支持一对多服务请
求传播，而不修改 NFD。

\subsection{应用特定 NDN 框架}

其他 NDN 系统聚焦特定领域。DNMP 使用 NDN 同步和层次 topic 构建安全网络
测量框架~\cite{dnmp}。MIA-NDN 探索 NDN 中以微服务为中心的 Interest 聚合
~\cite{miandn}。SECaaS 研究 multi-access edge 中的 security-as-a-service
~\cite{secaas}。这些项目表明，当名字和应用语义一起设计时，NDN 能够适配
实际服务型 workload。

限制在于通用性。应用特定框架常常假设特定数据类型、服务领域、聚合模式或
部署拓扑。NDNSF 则试图提供可复用框架层，支持 UAV 控制、视频、目标检测、
分布式推理、存储和协作等不同服务。拟实现的 UAV 和 DistributedInference
应用将检验该框架是否足够通用，而不仅仅适用于一个精心调优的应用。

\subsection{安全和访问控制机制}

NDN 默认对 Data packet 签名，但签名本身不能回答所有服务问题。服务框架
必须决定谁可以请求服务、谁可以提供服务、哪些数据应加密、密钥如何分发以
及重放消息如何被拒绝。属性基加密和 NAC-ABE 展示了细粒度访问控制如何与
NDN 数据名集成~\cite{goyal2006abe,dulal2023nacabe}。这些思想是 NDNSF 的核心。

开放问题是与服务工作流的集成。如果访问控制被临时添加到每个应用中，每个
应用都必须分别处理权限获取、加密属性、消息验证和防重放。NDNSF 将这些机
制纳入运行时：ServiceController 分发权限，服务消息根据服务级属性保护，
一次性 user/provider token 将 ACK、selection、Targeted fast path 和
response 绑定到当前调用。

\section{NDN 服务系统}

NDN Service Call（NSC）~\cite{nsc} 和其他 NDN 应用框架表明，NDN 能够在
没有 IP 端点绑定的情况下支持服务调用。NDN 的 Interest/Data 交换、基于名
字的路由、网内缓存和数据签名解决了 IP 系统中的若干问题。因此，这些系统
证明 NDN 是服务的有前景基础。

剩余缺口是：现有 NDN 服务系统尚未提供一个通用运行时来同时结合：

\begin{itemize}
  \item 通用动态服务 API；
  \item controller 授权的用户和提供者权限；
  \item 一对多请求发布；
  \item 服务提供者 ACK 元数据；
  \item 应用定义的选择策略；
  \item 自适应 admission control；
  \item 大数据引用；
  \item 服务提供者协作；
  \item 可在真实场景中评估的 service container。
\end{itemize}

\section{设计模式和权衡}

相关工作显示出几个反复出现的设计模式。首先，服务框架必须将服务名与提供
者位置解耦。第二，大型输入和输出应表示为命名数据对象，而不是直接嵌入小
请求 payload。第三，服务状态通常不能只依赖 NDN Pending Interest Table，
因为服务执行可能长时间、多步骤或协作式。第四，安全性必须成为调用工作流
的一部分，而不是事后附加。

这些模式也暴露出权衡。函数中心名字表达力强，但部署困难；serverless 放置
灵活，但需要编排；RPC 式调用易于理解，但倾向一对一执行；broker 系统在时
间上解耦，但可能引入基础设施依赖。NDNSF 的设计点是将类似 RPC 的开发者
API 与 NDN 的数据中心名字、签名 Data、基于 Sync 的发布和应用可见的提供
者选择结合起来。

\section{开放问题}

表~\ref{tab:framework-gaps-ch} 总结了推动 NDNSF 的开放问题。

\begin{table}[H]
\centering
\caption{现有面向服务框架中的开放问题。}
\label{tab:framework-gaps-ch}
\begin{tabularx}{\textwidth}{p{0.25\textwidth}X}
\toprule
问题 & 重要性 \\
\midrule
端点绑定 & 移动服务和间歇连接提供者应按名字和能力调用，而不是依赖稳定主机地址。 \\
一对多调用 & 可靠性和协作常常需要多个提供者听到请求，并允许用户选择一个或多个提供者。 \\
数据中心安全 & 数据在静态存储时也应保持签名、可验证和可选加密，而不只是经过安全信道时受保护。 \\
提供者授权 & 用户不仅要验证自己是否能调用服务，也要验证被选择的提供者是否有权提供服务。 \\
提供者选择 & 负载、队列长度、位置和可用性等元数据应成为服务流程的一部分。 \\
大数据 & 模型、视频、图片、日志和中间张量不应被迫塞进小请求 payload。 \\
协作 & 任务可能需要多个提供者、补偿请求或部分结果聚合。 \\
操作验证 & 真实系统需要可重复的证书、路由、配置和实验环境，以便在 toy
example 之外验证框架。 \\
\bottomrule
\end{tabularx}
\end{table}

NDNSF 旨在使用 NDN 名字、可验证 Data、基于 Sync 的发布和服务级授权来解
决这些问题。

本章通过把 related work 转化为 framework requirements 来支持 RQ1，而不
只是列出已有系统。它也为 RQ2 和 RQ3 铺垫：哪些性质必须由 NDNSF runtime
提供，而不能让每个 UAV 或推理应用分别重新实现。
